% Section 3 - Smart pointers
% Roberto Masocco <roberto.masocco@uniroma2.it>
% April 22, 2026

% ### Smart pointers ###
\section{Smart pointers}
\graphicspath{{figs/section3/}}

% --- Smart pointers ---
\begin{frame}{Smart pointers}{The problem with raw pointers}
	\justifying
	Raw pointers (\texttt{T *}) give full control over heap memory, but place the entire \textbg{burden of correctness} on the programmer:
	\begin{itemize}
		\item every \texttt{new} must be \textbg{matched} by exactly one \texttt{delete};
		\item in the presence of \textbg{multiple exit paths}, exceptions, or callbacks, this matching becomes error-prone;
		\item \textbg{ownership} — who is responsible for deleting — is implicit and easy to lose track of.
	\end{itemize}
	\begin{alertblock}{In a callback-driven, multithreaded system}
		\justifying
		An object created in one thread, passed to a callback, and deleted somewhere else is a recipe for \textbr{use-after-free bugs} that are nearly impossible to reproduce and debug.
	\end{alertblock}
	\textbg{Smart pointers} apply \textbg{RAII to heap memory}: the pointer object owns the resource and releases it automatically when destroyed.
\end{frame}
\begin{frame}{Smart pointers}{\texttt{std::unique\_ptr}}
	\justifying
	\texttt{std::unique\_ptr<T>} expresses \textbg{sole, exclusive ownership} of a heap-allocated object.\\
  \bigskip
	There is \textbg{exactly one} \texttt{unique\_ptr} per object, the object is destroyed when the pointer goes out of scope.\\
  \bigskip
	This enforces at compile time that no two owners exist simultaneously.\\
  \bigskip
	Use it as the \textbg{default} choice when ownership is clear and singular.\\
  \bigskip
  It cannot be \textbg{copied}, only \textbg{moved} (see code example).
\end{frame}
\begin{frame}[fragile]{Smart pointers}{\texttt{std::unique\_ptr}}
  \begin{columns}
		\column{.9\textwidth}
		\begin{lstlisting}[style=codebeamer, language=C++, caption={\texttt{unique\_ptr} creation, transfer, and automatic release}.]
#include <memory>

{
  auto p = std::make_unique<int>(42); // preferred creation
  std::cout << * p << std::endl;      // 42

  auto q = std::move(p); // ownership transferred to q
                         // p is now empty (nullptr)
} // q goes out of scope: int is deleted automatically
    \end{lstlisting}
	\end{columns}
\end{frame}
\begin{frame}[fragile]{Smart pointers}{\texttt{std::shared\_ptr}}
	\justifying
	\texttt{std::shared\_ptr<T>} allows \textbg{shared ownership}: multiple pointers may refer to the same object.\\
  \bigskip
	It maintains a \textbg{reference count}: the object is \textbg{destroyed} only when the \textbg{last} \texttt{shared\_ptr} to it is destroyed.\\
  \bigskip
	Always prefer \texttt{std::make\_shared} over \texttt{new}: it allocates the \textbg{object} and the \textbg{control block} in a \textbg{single allocation}, and is \textbg{exception-safe}.
\end{frame}
\begin{frame}[fragile]{Smart pointers}{\texttt{std::shared\_ptr}}
  \begin{columns}
		\column{.9\textwidth}
		\begin{lstlisting}[style=codebeamer, language=C++, caption=Reference counting with \texttt{shared\_ptr}.]
#include <memory>

auto a = std::make_shared<int>(7); // count = 1
{
  auto b = a;                      // count = 2
  std::cout << a.use_count() << std::endl; // 2
  auto c = a;                      // count = 3
} // b and c destroyed: count drops to 1

// a still valid; int alive
std::cout << a.use_count() << std::endl;   // 1
// a goes out of scope: count = 0, int destroyed
    \end{lstlisting}
	\end{columns}
\end{frame}
\begin{frame}{Smart pointers}{\texttt{std::weak\_ptr}}
	\justifying
	\texttt{std::weak\_ptr<T>} is a \textbg{non-owning observer}: it refers to an object managed by a \texttt{shared\_ptr} without incrementing the reference count.\\
	\bigskip
	It is used to \textbg{break ownership cycles} — two objects each holding a \texttt{shared\_ptr} to the other would keep each other alive forever.\\
	\bigskip
	To access the object, a \texttt{weak\_ptr} must be \textbg{promoted} to a \texttt{shared\_ptr} first, which \textbg{may fail} if the object has already been destroyed.
\end{frame}
\begin{frame}[fragile]{Smart pointers}{\texttt{std::weak\_ptr}}
	\begin{columns}
		\column{.9\textwidth}
		\begin{lstlisting}[style=codebeamer, language=C++, caption=Promoting a \texttt{weak\_ptr} to check if the object is still alive.]
std::shared_ptr<int> owner = std::make_shared<int>(99);
std::weak_ptr<int>   observer = owner;

if (auto p = observer.lock())   // returns shared_ptr or null
  std::cout << *p << std::endl; // object still alive
    \end{lstlisting}
	\end{columns}
\end{frame}
\begin{frame}[fragile]{Smart pointers}{The \texttt{SharedPtr} alias}
	\justifying
	In large codebases, writing \texttt{std::shared\_ptr<MyClass>} everywhere is cumbersome.\\
	A common C++ convention — and a \textbg{systematic pattern in ROS 2} — is to define a \texttt{SharedPtr} alias inside the class itself.
	\begin{columns}
		\column{.9\textwidth}
		\begin{lstlisting}[style=codebeamer, language=C++, caption=The \texttt{SharedPtr} alias as used in ROS 2 classes.]
class MyNode
{
public:
  using SharedPtr = std::shared_ptr<MyNode>;
  // ...
};
    \end{lstlisting}
	\end{columns}
	You will encounter \texttt{Node::SharedPtr}, \texttt{Publisher::SharedPtr}, \texttt{Subscription::SharedPtr} and many others as soon as you open any rclcpp example.\\
	They are all \texttt{std::shared\_ptr} instantiations behind the alias.
\end{frame}
